% concatenated from NS-convolution.tex
\documentclass{amsart}
\usepackage{enumerate}
\usepackage{amssymb}
\usepackage{mathtools}
\usepackage{amsmath}
\usepackage{color}
\usepackage[
colorlinks=true,
linkcolor=blue,
citecolor=blue,
urlcolor=blue]{hyperref}

\numberwithin{equation}{section}

\newcommand{\normmm}[1]{{\left\vert\kern-0.25ex\left\vert\kern-0.25ex\left\vert #1
\right\vert\kern-0.25ex\right\vert\kern-0.25ex\right\vert}}

\newtheorem{thm}{Theorem}[section]
\newtheorem{lem}[thm]{Lemma}
\newtheorem{defi}[thm]{Definition}
\newtheorem{rem}[thm]{Remark}
\newtheorem{cor}[thm]{Corollary}
\newtheorem{nota}[thm]{Notation}
\newtheorem{prop}[thm]{Proposition}
\newtheorem{fact}[thm]{Fact}
\newtheorem{exmaple}[thm]{Example}
\newtheorem{prob}[thm]{Problem}
\newtheorem{conj}[thm]{Conjecture}
\newtheorem{question}[thm]{Question}

\newtheorem{thmx}{Theorem}[section]
\renewcommand{\thethmx}{\Alph{thmx}}
\newtheorem{corx}[thmx]{Corollary}
\renewcommand{\thecorx}{\Alph{corx}}



\newcommand{\E}{{\mathcal{E}}}
\newcommand{\N}{{\mathbb N}}
\newcommand{\R}{{\mathcal{R}}}
\newcommand{\Z}{{\mathbb{Z}^{n}_{m}}}
\newcommand{\dc}{\mathbf{d}}
\newcommand{\La}{\mathbf{\mathcal{L}}}
\newcommand{\tr}{\mathrm{tr}}
\newcommand{\Fr}{\mathfrak{F}}
\newcommand{\F}{\mathcal{F}}
\newcommand{\M}{\mathbb{M}}
\newcommand{\al}{\alpha}
\newcommand{\s}{\mathbf{s}}
\newcommand{\Ex}{\mathcal{E}}
\newcommand{\ran}{\mathrm{range}}
\newcommand{\Wk}{\overline{\mathcal{W}_{k}}^{\|\cdot\|_{H_{1}}}}
\newcommand{\Vk}{\overline{\mathcal{V}_{k}}^{\|\cdot\|_{H_{1}}}}
\newcommand{\supp}{{\mathrm{supp}}}
\newcommand{\var}{{\mathrm{var}}}

\begin{document}


\title[Strong Convolution Inequality]{The Strong Naor--Schechtman Convolution Inequality on the Product of Cyclic Groups}


\author[Jiao]{Yong Jiao}
\address{School of Mathematics and Statistics, HNP-LAMA, Central South University, Changsha 410075, China}
\email{jiaoyong@csu.edu.cn}


\author[Luo]{Sijie Luo}
\address{School of Mathematics and Statistics, HNP-LAMA, Central South University, Changsha 410075, China}
\email{sijieluo@csu.edu.cn}

\author[Zanin]{Dmitriy Zanin}
\address{School of Mathematics and Statistics, HNP-LAMA, Central South University, Changsha 410075, China}
\email{d.zanin@unsw.edu.au}

\author[Zhou]{Dejian Zhou}
\address{School of Mathematics and Statistics, HNP-LAMA, Central South University, Changsha 410075, China}
\email{zhoudejian@csu.edu.cn}
\date{}

\begin{abstract}
In this paper, we use martingale methods to resolve the convolution inequality problem posed by Naor and Schechtman in \cite[Question 6.1]{N-S2016} (see also \cite{Na2016}). More precisely, we establish the following strong convolution inequality on products of finite cyclic groups. For each $1<p<\infty$ and every $n$, $m\in \mathbb{N}$ we have
\[
\begin{split}
&\sum_{\varepsilon\in \{-1,1\}^{n}}\sum_{x\in \Z}\left|E_{\{1,\cdots,n\}}f(x+\varepsilon)-E_{\{1,\cdots,n\}}f(x-\varepsilon)\right|^{p}\\
\leq& (p^{*}-1)^{p}\sum_{\varepsilon\in \{-1,1\}^{n}}\sum_{x\in \Z}\left|\varepsilon_{j}\left[E_{\{1,\cdots,n\}\setminus\{j\}}f(x+e_{j})-E_{\{1,\cdots,n\}\setminus\{j\}}f(x-e_{j})\right]\right|^{p},
\end{split}
\]
where $p^{*}=\max\{p,p/(p-1)\}$ and $E_{A}f(x)=\frac{1}{2^{n}}\sum_{\delta\in \{-1,1\}^{n}}f\left(x+\sum_{j\in A}\delta_{j}e_{j}\right)$ for every $A\subseteq \{1,\cdots,n\}$.
\end{abstract}

\subjclass[2020]{Primary 46B20; Secondary 60G42, 42C10.}
\keywords{Convolution inequality, martingales, metric $X_{p}$ inequality}
\maketitle

\section{Introduction}

The study of metric property determined only by the distance of Banach spaces forms a central theme in the Ribe program. The metric $X_{p}$ inequality introduced by Naor and Schechtman \cite{N-S2016} is a nonlinear counterpart of the Rosenthal inequality, which serves as an obstruction to certify the failure of bi-Lipschitz emebddability of $L_{q}$ into $L_{p}$ for $2<q<p<\infty$.

To make statements precisely, we introduce necessary notations that will be used throughly. For a given $n$, let $\Z$ be the $n$-folds Cartesian product of the cyclic group $\mathbb{Z}_{m}$. We denote that $\Omega_{n}=\{-1,1\}^{n}$ the $n$-dimensional hypercube. Both $\Z$ and $\Omega_{n}$ are equipped with the uniform probability measures $\nu$ and $\mu$, respectively. Let $[n]=\{1,\cdots, n\}$, and for each $j\in [n]$ we define the translation operator $T_{j}$ by
\[
T_{j}f(x)\coloneqq f(x+e_{j}),\qquad f:\Z\to \mathbb{R},
\]
and $e_{j}=(0,\cdots, 1,\cdots,0)$ is the vector with $1$ appears in the $j$-th position. It is clear that the linear operator $T_{j}$ is an isometry on $L_{p}(\Z)$ for all $p\in [1,\infty]$ with its inverse given by
\[
T^{-1}_{j}f(x)=f(x-e_{j}),\qquad f:\Z\to \mathbb{R}.
\]
%It is easy to verify that for each $j\in [n]$,
%\[
%\|T_{j}f\|^{p}_{L_{p}(\Z)}=\frac{1}{m^{n}}\sum_{x\in \Z}|T_{j}f(x)|^{p}=\frac{1}{m^{n}}\sum_{x\in \Z}|f(x+e_{j})|^{p}=\|f\|^{p}_{L_{p}(\Z)},\quad 1\leq p<\infty,
%\]
%and
%\[
%\|T_{j}f\|_{L_{\infty}(\Z)}=\max_{x\in \Z}|T_{j}f(x)|=\max_{x\in \Z}|f(x+e_{j})|=\max_{x\in \Z}|f(x)|.
%\]
The $j$-th expectation operator $E_{j}$ and the $j$-th derivative operator $D_{j}$ are defined as follows
\[
E_{j}\coloneqq \frac{T_{j}+T^{-1}_{j}}{2}\quad\mbox{and}\quad D_{j}=\frac{T_{j}-T^{-1}_{j}}{2}.
\]
It is clear that the family $\{E_{j}\}_{j=1}^{n}$ commute with each other. For each $A\subseteq [n]$, let $E_{A}\coloneqq \prod_{j\in A}E_{j}$ and $E_{\emptyset}=\mathrm{Id}$. Hence,
\[
E_{A}f(x)=\frac{1}{2^{|A|}}\sum_{\delta\in \{-1,1\}^{A}}f\left(x+\sum_{j\in A}\delta_{j}e_{j}\right)=\frac{1}{2^{n}}\sum_{\delta\in \Omega_{n}}f\left(x+\sum_{j\in A}\delta_{j}e_{j}\right),
\]
for every $f:\Z\to \mathbb{R}$, where $|A|$ stands for the cardinality of $A$. For $A\subseteq [n]$ and $\varepsilon\in \Omega_{n}$, we let
\[
T^{\varepsilon}_{A}\coloneqq \prod_{j\in A}T^{\varepsilon_{j}}_{j}.
\]



The metric $X_{p}$ inequality of Naor and Schechtman is stated as follows, which is one of the essential result in \cite{N-S2016}.
\begin{thm}[Naor--Schechtman]
For $2\leq p<\infty$, and $k$, $m\in \mathbb{N}$ with $k\in [n]$ and
\[
m\geq \frac{n^{3/2}\log(p)}{\sqrt{k}}+pn,
\]
then there exists a constant $c_{p}>0$ depending only on $p$ such that the following inequality holds for every $f:\mathbb{Z}^{n}_{4m}\to L_{p}$
\[
\begin{split}
&\frac{(p\log(p))^{-1}}{\binom{n}{k}}\sum_{\substack{S\subseteq[n] \\|S|=k}}\frac{\mathbb{E}\left[\left\|f(x+2m\varepsilon_{S})-f(x)\right\|^{p}_{L_{p}}\right]}{m^{p}}\\
\leq& c_{p}\left(\frac{k}{n}\sum_{j=1}^{n}\mathbb{E}\left\|T_{j}f(x)-f(x)\right\|^{p}_{L_{p}}+\left(\frac{k}{n}\right)^{p/2}\mathbb{E}\left\|T^{\varepsilon}_{[n]}f(x)-f(x)\right\|^{p}_{L_{p}}\right),
\end{split}
\]
where the expectation is with respect to $(x,\varepsilon)\in \mathbb{Z}^{n}_{4m}\times \{-1,1\}^{n}$ chosen unifromly at random.
\end{thm}
\noindent In the same paper, Naor and Schechtman showed that the scaling appears in the metric $X_{p}$ inequality is necessary (see \cite[Propostion 1.4]{N-S2016}). This leads them to conjecture that $L_{p}$ space is an $X_{p}$ metric space with optimal scaling parameter.
\begin{conj}[Naor--Schechtman \cite{N-S2016}]\label{N-S conjecture}
For every $2<p<\infty$ there exist positive constants $\alpha_{p}$ and $c_{p}$ such that for every $n$, $m\in \mathbb{N}$ with $k\in[n]$, $m\geq C_{p}\sqrt{n/k}$, then for every $f:\mathbb{Z}^{n}_{4m}\to \mathbb{R}$ we have
\[
\begin{split}
&\frac{\alpha_{p}}{\binom{n}{k}}\sum_{\substack{S\subseteq [n] \\|S|=k}}\frac{\mathbb{E}\left[\left|f(x+2m\varepsilon_{S})-f(x)\right|^{p}\right]}{m^{p}}\\
\leq& c_{p}\left(\frac{k}{n}\sum_{j=1}^{n}\left\|T_{j}f-f\right\|_{L_{p}(\Z\times \Omega_{n})}+\left(\frac{k}{n}\right)^{p/2}\mathbb{E}\left\|T^{\varepsilon}_{[n]}f-f\right\|^{p}_{L_{p}(\Z\times \Omega_{n})}\right),
\end{split}
\]
where the expectation is with respect to $(x,\varepsilon)\in \mathbb{Z}^{n}_{4m}\times \{-1,1\}^{n}$ chosen unifromly at random.
\end{conj}

In \cite[Section 6]{N-S2016}, Naor and Schechtman proposed an approach to Conjecture \ref{N-S conjecture} based on a convolution inequality, showing
that the inequality formulated in \cite[Question 6.1]{N-S2016} would imply the conjecture. This convolution inequality therefore provides an analytic route to the optimal scaling and its consequences for the nonlinear geometry of $L_{p}$ spaces. Subsequently, Naor \cite{Na2016} resolved the scaling conjecture using Lust-Piquard's bounds for discrete Riesz transforms \cite{LP1998},
while explicitly noting that his argument left Question $6.1$ unresolved. More recently, Areshidze \cite{Are2026} obtained the optimal $p/\log p$ bounds for higher-order Rademacher chaos via a martingale argument. Related developments have also been made in the study of discrete Riesz transform estimates. Domelevo, Ivanisvili, Petermichl, and Volberg \cite{D-I-P-V2026} gave a commutative Bellman-function proof of the discrete Riesz transform bound for $2\leq p<\infty$, with a constant bounded by $C(p-1)$, where $C$ is universal. In the complementary range $1<p<2$, Jiao, Luo, Zanin, and Zhou \cite{JLZZ2026} introduced an exponential lifting technique and
combined it with martingale methods to establish the sharp fractional Riesz estimate, while Xu and Zhang \cite{X-Z2026} independently established the same estimate using noncommutative BMO theory, thereby resolving the problem posed by Efraim and Lust-Piquard \cite{ELP2008}.

To state \cite[Question 6.1]{N-S2016}, we introduce two basic mappings. For $j\in [n]$ and every $f:\Z\to \mathbb{R}$ and $j\in [n]$, we define functions $H_{f}: \Z\times \Omega_{n}\to \mathbb{R}$ by
\[
H_{f}(x,\varepsilon)=(E_{[n]} f)(x+\varepsilon)-(E_{[n]} f)(x-\varepsilon),
\]
\[
b^{(j)}_{f}(x)\coloneqq E_{[n]\setminus\{j\}}f(x+e_{j})-E_{[n]\setminus \{j\}}f(x-e_{j})=2D_{j}E_{[n]\setminus\{j\}}(f)(x).
\]
Question $6.1$ of \cite{N-S2016} asks whether there exists a constant $c_{p}>0$ such that for any $n$, $m\in \mathbb{N}$ the following inequality holds for every $f:\Z\to \mathbb{R}$
\begin{equation}\label{concolution inequality of N-S}
\|H_{f}\|_{L_{p}(\Z\times\Omega_{n})}\leq c_{p}\left(\left\|\sum_{j=1}^{n}\varepsilon_{j}b^{(j)}_{f}\right\|_{L_{p}(\Z\times \Omega_{n})}+\sum_{j=1}^{n}\left\|T_{j}f-f\right\|_{L_{p}(\Z\times \Omega_{n})}\right).
\end{equation}

In this paper, we settle the problem of Naor and Schechtman \cite[Question 6.1]{N-S2016} by establishing the following strong convolution inequality.

\begin{thm}\label{NS-main}
For every $1<p<\infty$ and every $n$, $m\in \N$, the following inequality holds for every $f:\Z\to \mathbb{R}$
\begin{equation}
\left\|H_{f}\right\|_{L_{p}(\Z\times \Omega_{n})}\leq (p^{*}-1)\left\|\sum_{j=1}^{n}\varepsilon_{j}b^{(j)}_{f}\right\|_{L_{p}(\Z\times \Omega_{n})},
\end{equation}
where $p^{*}=\max\{p,\frac{p}{p-1}\}$.
\end{thm}
\noindent The rest of the paper is devoted to providing a proof of Theorem \ref{NS-main}.

\section{The proof of the main result}
This section is devoted to establishing Theorem \ref{NS-main}, and we begin with introducing further notations. For each $A\subseteq [n]$ and $\varepsilon\in \Omega_{n}$,  let  
\[
S^{\varepsilon}_{A}\coloneqq \frac{T^{\varepsilon}_{A}+T^{-\varepsilon}_{A}}{2}.
\]
For a given $f:\Z\to \mathbb{R}$, we define
\[
\mathcal{S}f(x,\varepsilon)\coloneqq \frac{f(x+\varepsilon)+f(x-\varepsilon)}{2}=\frac{T^{\varepsilon}_{[n]}f(x)+T^{-\varepsilon}_{[n]}f(x)}{2}.
\]
In the sequel, for each $j\in[n]$, we let $\varepsilon_{j}:\{-1,1\}^{n}\to \{-1,1\}$ be the $j$-th coordinate function. Let $\mathcal{B}$ be the Borel $\sigma$-algebra on $\Z$. We define a sequence of $\sigma$-algebras on $\Z\times \Omega_{n}$ by setting
\[
\mathcal{F}_0=\{B\times \Omega_{n}: B\in\mathcal{B}\}
\]
and
\[
\mathcal{F}_{j}\coloneqq \sigma\{B\times A: B\in \mathcal{B}, A\subseteq \Omega_{j}\}\qquad  j\in [n].
\]
Then, $\mathcal{F}_{0}\subseteq \mathcal{F}_{1}\subseteq \cdots\subseteq \mathcal{F}_{n}$ forms a filtration on the probability space $\Z\times \Omega_{n}$.

The following lemma provides a representation of $\mathcal{S}(b^{(j)}_{f})$ via the conditional expectation.
\begin{lem}\label{representation 1}
For every $f:\Z\to \mathbb{R}$ and  each $j\in [n]$, the following identity holds
\[
\begin{split}
\mathbb{E}\left[\mathcal{S}b^{(j)}_{f}\Big{|}\mathcal{F}_{j-1}\right](x,\varepsilon)=E_{[n]\setminus[j-1]}S^{\varepsilon}_{[j-1]}b^{(j)}_{f}(x),\quad (x,\varepsilon)\in \Z\times \Omega_{n}.
\end{split}
\]
\end{lem}
\begin{proof}
By the definition of $\mathcal{F}_{j-1}$, we have 
\begin{equation}\label{expansion 1}
\begin{split}
&\mathbb{E}\left[\mathcal{S}b^{(j)}_{f}\Big{|}\mathcal{F}_{j-1}\right](x,\varepsilon)\\
=&\frac{1}{2^{n-j+1}}\sum_{\varepsilon^{\prime}\in \Omega_{n-j+1}}\frac{b^{(j)}_{f}(x+\sum_{l=1}^{j-1}\varepsilon_{l}e_{l}+\sum_{l=j}^{n}\varepsilon^{\prime}_{l-j-1}e_{l})+b^{(j)}_{f}(x-\sum_{l=1}^{j-1}\varepsilon_{l}e_{l}-\sum_{l=j}^{n}\varepsilon^{\prime}_{l-j-1}e_{l})}{2}\\
=&\frac{1}{2^{n-j+1}}\sum_{\varepsilon^{\prime}\in \Omega_{n-j+1}}\frac{\left(T^{(\varepsilon_{1},\cdots, \varepsilon_{j-1},\varepsilon^{\prime}_{1},\cdots, \varepsilon^{\prime}_{n-j+1})}_{[n]}b^{(j)}_{f}\right)(x)+\left(T^{(-\varepsilon_{1},\cdots, -\varepsilon_{j-1},-\varepsilon^{\prime}_{1},\cdots, -\varepsilon^{\prime}_{n-j+1})}_{[n]}b^{(j)}_{f}\right)(x)}{2}\\
=&D_{j}E_{[n]\setminus\{j\}}\left(\frac{1}{2^{n-j+1}}\sum_{\varepsilon^{\prime}\in \Omega_{n-j+1}}\left(T^{(\varepsilon_{1},\cdots, \varepsilon_{j-1},\varepsilon^{\prime}_{1},\cdots, \varepsilon^{\prime}_{n-j+1})}_{[n]}b^{(j)}_{f}\right)(x)\right)\\
&+D_{j}E_{[n]\setminus\{j\}}\left(\frac{1}{2^{n-j+1}}\sum_{\varepsilon^{\prime}\in \Omega_{n-j+1}}\left(T^{(-\varepsilon_{1},\cdots, -\varepsilon_{j-1},-\varepsilon^{\prime}_{1},\cdots, -\varepsilon^{\prime}_{n-j+1})}_{[n]}b^{(j)}_{f}\right)(x)\right).
\end{split}
\end{equation}
It is clear that
\begin{equation}\label{expansion 2}
\frac{1}{2^{n-j+1}}\sum_{\varepsilon^{\prime}\in \Omega_{n-j+1}}\left(T^{(\varepsilon_{1},\cdots, \varepsilon_{j-1},\varepsilon^{\prime}_{1},\cdots, \varepsilon^{\prime}_{n-j+1})}_{[n]}b^{(j)}_{f}\right)(x)=T^{\varepsilon}_{[j-1]}E_{[n]\setminus[j-1]}f(x),
\end{equation}
and
\begin{equation}\label{expansion 3}
\frac{1}{2^{n-j+1}}\sum_{\varepsilon^{\prime}\in \Omega_{n-j+1}}\left(T^{(-\varepsilon_{1},\cdots, -\varepsilon_{j-1},-\varepsilon^{\prime}_{1},\cdots, -\varepsilon^{\prime}_{n-j+1})}_{[n]}b^{(j)}_{f}\right)(x)=T^{-\varepsilon}_{[j-1]}E_{[n]\setminus[j-1]}f(x).
\end{equation}
Substituting \eqref{expansion 2} and \eqref{expansion 3} into \eqref{expansion 1} yields that
\[
\mathbb{E}\left[\mathcal{S}b^{(j)}_{f}\Big{|}\mathcal{F}_{j-1}\right](x,\varepsilon)=E_{[n]\setminus[j-1]}S^{\varepsilon}_{[j-1]}b^{(j)}_{f}(x),
\]
where we used the fact that all operators are commute with each other and the fact that $b^{(j)}_{f}=2D_{j}E_{[n]\setminus\{j\}}f$.
\end{proof}

Applying Lemma \ref{representation 1}, we obtain the following decomposition of $H_{f}$, which is one of the key ingredients in the proof of our main result.
\begin{lem}\label{main decomposition}
For every $f:\Z\to \mathbb{R}$, we have
\begin{equation}\label{decomposition of Hf}
H_{f}(x,\varepsilon)=\sum_{j=1}^{n}\varepsilon_{j}\mathbb{E}\left[\mathcal{S}b^{(j)}_{f}\Big{|}\mathcal{F}_{j-1}\right](x,\varepsilon), \quad (x,\varepsilon)\in \Z\times \Omega_{n}.
\end{equation}
\end{lem}
\begin{proof}
Note here that $H_{f}=E_{[n]}\left(T^{\varepsilon}_{[n]}-T^{-\varepsilon}_{[n]}\right)f$. For each $j\in [n]$, we define
\begin{equation}\label{decomposition 1}
\mathcal{M}_{j}(f)\coloneqq E_{[n]}E_{[n]\setminus[j]}\left(T^{\varepsilon}_{[j]}-T^{-\varepsilon}_{[j]}\right)(f),
\end{equation}
and let $\mathcal{M}_{0}(f)\equiv 0$ for convenience. Since $E_{j}=\frac{T_{j}+T^{-1}_{j}}{2}$ and $D_{j}=\frac{T_{j}-T^{-1}_{j}}{2}$, it is clear that
\begin{equation}\label{key 1}
T_{j}=E_{j}+D_{j}\qquad\mbox{and}\qquad T^{-1}_{j}=E_{j}-D_{j}.
\end{equation}
By $T^{\varepsilon}_{[j]}=T^{\varepsilon_{j}}_{j}T^{\varepsilon}_{[j-1]}$ and $T^{-\varepsilon}_{[j]}=T^{-\varepsilon_{j}}_{j}T^{-\varepsilon}_{[j-1]}$, it follows from \eqref{key 1} that
\begin{equation}\label{decomposition 2}
T^{\varepsilon}_{[j]}-T^{-\varepsilon}_{[j]}=(E_{j}+\varepsilon_{j}D_{j})T^{\varepsilon}_{[j-1]}-(E_{j}-\varepsilon_{j}D_{j})T^{-\varepsilon}_{[j-1]}.
\end{equation}
Substituting \eqref{decomposition 2} into \eqref{decomposition 1} yields that
\begin{equation}\label{decomposition 3}
\begin{split}
\mathcal{M}_{j}(f)&=E_{[n]}E_{[n]\setminus[j]}E_{j}\left(T^{\varepsilon}_{[j-1]}-T^{-\varepsilon}_{[j-1]}\right)(f)+2\varepsilon_{j}E_{[n]}E_{[n]\setminus[j]}D_{j}S^{\varepsilon}_{[j-1]}(f)\\
&=\mathcal{M}_{j-1}(f)+2\varepsilon_{j}E_{[n]}E_{[n]\setminus[j]}D_{j}S^{\varepsilon}_{[j-1]}(f)\\
&=\mathcal{M}_{j-1}(f)+\varepsilon_{j}E_{[n]\setminus[j-1]}S^{\varepsilon}_{[j-1]}b^{(j)}_{f}\\
&=\mathcal{M}_{j-1}(f)+\varepsilon_{j}\mathbb{E}\left[\mathcal{S}b^{(j)}_{f}|\mathcal{F}_{j-1}\right],
\end{split}
\end{equation}
where the last equality is due to Lemma  \ref{representation 1}. Noting $M_{n}(f)=H_{f}$,  by induction of \eqref{decomposition 3},  we obtain
\[
H_{f}=\sum_{j=1}^{n}\varepsilon_{j} \mathbb{E}\left[\mathcal{S}b^{(j)}_{f}\Big{|}\mathcal{F}_{j-1}\right].
\]
\end{proof}
The next elementary lemma provides a connection between the martingale difference operator and Lemma \ref{main decomposition}.
\begin{lem}\label{representation of martingale difference}
For each $f:\Z\times \Omega_{n}\to \mathbb{R}$, we have
\[
d_{j}(f)\coloneqq\mathbb{E}\left[f\Big{|}\mathcal{F}_{j}\right]-\mathbb{E}\left[f\Big{|}\mathcal{F}_{j-1}\right]=\varepsilon_{j}\mathbb{E}\left[\varepsilon_{j}f\Big{|}\mathcal{F}_{j-1}\right],\quad1\leq j\leq n.
\]
\end{lem}
\begin{proof}
The proof follows from the Walsh expansion directly. Indeed, for $f:\Z\times \Omega_{n}\to \mathbb{R}$, we define $F:\Z\to L_{2}(\Omega_{n};\mathbb{R})$ by
\[
F(x)\coloneqq f(x,\cdot),\quad x\in \Z.
\]
Since $\{w_{A}\}_{A\subseteq [n]}$ forms an orthogonal basis of $L_{2}(\Omega_{n})$, it follows that for each $x\in \Z$,  we have
\[
F(x)=\sum_{A\subseteq [n]}\widehat{F(x)}(A)w_{A}.
\]
By the definition of $\{\mathcal{F}_{j}\}_{j=0}^{n}$, we have
\[
\mathbb{E}[F|\mathcal{F}_{j}](x)=\sum_{A\subseteq [j]}\widehat{F(x)}(A)w_{A}.
\]
Thus,
\begin{equation}\label{Walsh expansion 1}
\begin{split}
d_{j}(f)(x,\cdot)&=\sum_{\substack{A\subseteq [n]\\\max(A)=j}}\widehat{F(x)}(A)w_{A}\\
&=\sum_{B\subseteq [j-1]}\widehat{F(x)}(B)w_{B\cup\{j\}}=\varepsilon_{j}\cdot\left(\sum_{B\subseteq [j-1]}\widehat{F(x)}(B)w_{B}\right).
\end{split}
\end{equation}
On th other hand, by the definition of conditional expectation, we have
\begin{equation}\label{Walsh expansion 2}
\begin{split}
\mathbb{E}\left[\varepsilon_{j} F\Big{|}\mathcal{F}_{j-1}\right](x,\cdot)&=\mathbb{E}\left[\sum_{j\in A\subseteq [n]}\widehat{F(x)}(A)w_{A\setminus\{j\}}\Big{|}\mathcal{F}_{j-1}\right]\\
&\quad +\mathbb{E}\left[\sum_{A\subseteq [n]\setminus\{j\}}\widehat{F(x)}(A)w_{A\cup\{j\}}\Big{|}\mathcal{F}_{j-1}\right]\\
&=\sum_{\substack{B\subseteq[n]\\\max(B)=j}}\widehat{F(x)}(B)w_{B\setminus\{j\}}=\sum_{B\subseteq [j-1]}\widehat{F(x)}w_{B}.
\end{split}
\end{equation}
Combining \eqref{Walsh expansion 1} with \eqref{Walsh expansion 2} yields the desired result.
\end{proof}

We now recall the \emph{sharp} inequality for martingale transform of Burkholder \cite{Bur1984}.
\begin{thm}[Burkholder]\label{Burkholder martingale transform}
Suppose that $M=\{M_{j}\}_{j=0}^{\infty}$ martingale with $M_{0}\equiv 0$. Let $d_{j}\coloneqq M_{j}-M_{j-1}$ be the martingale difference for each $j\in \mathbb{N}$. If $\{\varepsilon_{j}\}_{j=1}^{n}$ is a sequence of numbers in $\{-1,1\}$, then the following inequality holds
\[
\left\|\sum_{j=1}^{\infty}\varepsilon_{j}d_{j}\right\|_{L_{p}(\Omega)}\leq (p^{*}-1)\|M\|_{L_{p}(\Omega)},  \quad 1<p<\infty,
\]
where $p^{*}=\max\{p,p/(p-1)\}$
\end{thm}

We now provide the proof of Theorem \ref{NS-main}. 
\begin{proof}[Proof of Theorem \ref{NS-main}]
By Lemma \ref{main decomposition}, we have
\[
H_{f}(x,\varepsilon)=\sum_{j=1}^{n}\varepsilon_{j}\mathbb{E}\left[\mathcal{S}b^{(j)}_{f}|\mathcal{F}_{j-1}\right](x,\varepsilon).
\]
Hence, for each $g\in L_{p^{\prime}}(\Z\times \Omega_{n})$ with $\|g\|_{L_{p^{\prime}}}\leq 1$ and $p^{\prime}=\frac{p}{p-1}$, we have
\begin{equation}\label{martingale transform 1}
\begin{split}
\mathbb{E}_{\nu\otimes \mu}\left[gH_{f}\right]&=\mathbb{E}_{\nu\otimes \mu}\left[\left(\sum_{j=1}^{n}\mathbb{E}[\varepsilon_{j} g|\mathcal{F}_{j-1}]\right)\cdot\mathcal{S}b^{(j)}_{f}\right]\\
&=\mathbb{E}_{\nu\otimes \mu}\left[\sum_{j=1}^{n}d_{j}(g)\cdot \varepsilon_{j} \mathcal{S}b^{(j)}_{f}\right],
\end{split}
\end{equation}
where we used Lemma \ref{representation of martingale difference}.

Let $(\delta^{\prime}_{1},\cdots, \delta^{\prime}_{n})$ be an i.i.d. Rademacher sequence (on a probability space $(\Omega,\mathcal{B},\mu^{\prime})$), which is independent of $(\varepsilon_{1},\cdots,\varepsilon_{n})$. It is clear that for each $\{\alpha_{j}\}_{j=1}^{n}$ and $\{\beta_{j}\}_{j=1}^{n}\subseteq \mathbb{R}$ we have
\begin{equation}\label{martingale transform 2}
\sum_{j=1}^{n}\alpha_{j}\beta_{j}=\mathbb{E}_{\mu^{\prime}}\left[\left(\sum_{j=1}^{n}\delta_{j}\alpha_{j}\right)\cdot\left(\sum_{j=1}^{n}\delta_{j}\beta_{j}\right)\right].
\end{equation}
Substituting \eqref{martingale transform 2} into \eqref{martingale transform 1}  yields
\begin{equation}\label{martingale transform 3}
\begin{split}
|\mathbb{E}_{\nu\otimes \mu}\left[gH_{f}\right]|&=\mathbb{E}_{\nu\otimes \mu\otimes \mu^{\prime}}\left[\left(\sum_{j=1}^{n}\delta_{j}d_{j}(g)\right)\cdot \left(\sum_{j=1}^{n}\delta_{j}\varepsilon_{j} \mathcal{S}b^{(j)}_{f}\right)\right]\\
&\leq \left\|\sum_{j=1}^{n}\delta_{j}d_{j}(g)\right\|_{L_{p^{\prime}}(\Z\times \Omega_{n}\times \Omega)}\cdot\left\|\sum_{j=1}^{n}\delta_{j}\varepsilon_{j} \mathcal{S}b^{(j)}_{f}\right\|_{L_{p}(\Z\times \Omega_{n}\times \Omega)}
\end{split}
\end{equation}
Note that for each $j\in [n]$ the random variable $\delta_{j}\varepsilon_{j}$ has the same distribution of $\varepsilon_{j}$, then it follows that
\begin{equation}\label{martingale transform 4}
\left\|\sum_{j=1}^{n}\delta_{j}\varepsilon_{j} \mathcal{S}b^{(j)}_{f}\right\|_{L_{p}(\Z\times \Omega_{n}\times \Omega)}=\left\|\sum_{j=1}^{n}\varepsilon_{j}\mathcal{S}b^{(j)}_{f}\right\|_{L_{p}(\Z\times \Omega_{n})}\leq \left\|\sum_{j=1}^{n}\varepsilon_{j}b^{(j)}_{f}\right\|_{L_{p}(\Z\times \Omega_{n})},
\end{equation}
where we used the fact that $\mathcal{S}$ is a contraction.

We now apply Theorem \ref{Burkholder martingale transform} to bound $\left\|\sum_{j=1}^{n}\delta_{j}d_{j}(g)\right\|_{L_{p^{\prime}}(\Z\times \Omega_{n}\times \Omega)}$. For each $\omega\in \Omega$, by Theorem \ref{Burkholder martingale transform}, we have
\[
\left\|\sum_{j=1}^{n}\delta_{j}(\omega)d_{j}(g)\right\|_{L_{p^{\prime}}(\Z\times \Omega_{n})}\leq (p^{*}-1)\|g\|_{L_{p^{\prime}}(\Z\times \Omega_{n})}.
\]
By the Fubini theorem, it yields that
\begin{equation}\label{martingale transform 5}
\begin{split}
\left\|\sum_{j=1}^{n}\delta_{j}d_{j}(g)\right\|_{L_{p^{\prime}}(\Z\times \Omega_{n}\times \Omega)}&=\left(\mathbb{E}_{\mu^{\prime}}\left[\left\|\sum_{j=1}^{n}\delta_{j}(\omega)d_{j}(g)\right\|^{p^{\prime}}_{L_{p^{\prime}(\Z\times \Omega_{n})}}\right]^{p^{\prime}}\right)^{1/p^{\prime}}\\
&\leq (p^{*}-1)\|g\|_{L_{p^{\prime}}(\Z\times \Omega_{n})}\leq p^{*}-1.
\end{split}
\end{equation}
Combining \eqref{martingale transform 3}, \eqref{martingale transform 4}, and \eqref{martingale transform 5} together yields
\[
\|H_{f}\|_{L_{p}(\Z\times \Omega_{n})}=\sup_{\|g\|_{L_{p^{\prime}}(\Z\times \Omega_{n})}\leq 1}|\mathbb{E}_{\nu\otimes \mu}[gH_{f}]|\leq (p^{*}-1)\left\|\sum_{j=1}^{n}\varepsilon_{j}b^{(j)}_{f}\right\|_{L_{p}(\Z\times \Omega_{n})},
\]
which completes our proof.
\end{proof}


\begin{rem}
To bound the term $\left\|\sum_{j=1}^{n}\delta_{j}d_{j}(g)\right\|_{L_{p^{\prime}}(\Z\times \Omega_{n}\times \Omega)}$, we may use the Burkholder-Gundy inequality directly. However, the constant $c_{p}$ we obtain in this way will be worse than the sharp inequality of Burkholder as we shown as above.
\end{rem}

\section*{Acknowledgment}

This work was supported by the National Natural Science Foundation of China (Grant Nos. 12125109 \& W2411005 \& 12671165); the Natural Science Foundation of Hunan Province (Grant Nos: 2025ZYJ002, 2024JJ1010 \& 2024RC3040); the Scientific Research Fund of Hunan Provincial Education Department (Grant Nos. 25A0009 \& 25B0008).

During the preparation of this work, the authors used GPT-5.6 Sol to assist with language editing and refinement of mathematical arguments. The author takes full responsibility for the content of this paper.




%\bibliographystyle{amsplain}
%\bibliography{References}

\providecommand{\bysame}{\leavevmode\hbox to3em{\hrulefill}\thinspace}
\providecommand{\MR}{\relax\ifhmode\unskip\space\fi MR }
% \MRhref is called by the amsart/book/proc definition of \MR.
\providecommand{\MRhref}[2]{%
  \href{http://www.ams.org/mathscinet-getitem?mr=#1}{#2}
}
\providecommand{\href}[2]{#2}
\begin{thebibliography}{1}

\bibitem{Are2026}
N.~Areshidze, \emph{Sharp metric {$X_{p}$} inequalities via martingales}, 2026,
  Preprint, arXiv:2608.29367.

\bibitem{Bur1984}
D.~Burkholder, \emph{Boundary value problems and sharp inequalities for
  martingale transforms}, Ann. Probab. \textbf{12} (1984), no.~3, 647--702.
  \MR{744226}

\bibitem{D-I-P-V2026}
K.~Domelevo, P.~Ivanisvili, S.~Petermichl, and A.~Volberg, \emph{Dimension-free
  bounds for {R}iesz transforms on the {H}amming cube via a {B}ellman
  function}, 2026, Preprint, arXiv:2606.20289.

\bibitem{ELP2008}
L.~B. Efraim and F.~Lust-Piquard, \emph{Poincar\'{e} type inequalities on the
  discrete cube and in the {CAR} algebra}, Probab. Theory Related Fields
  \textbf{141} (2008), no.~3-4, 569--602.

\bibitem{JLZZ2026}
Y.~Jiao, S.~Luo, D.~Zanin, and D.~Zhou, \emph{Sharp fractional {R}iesz
  estimates on the hypercube}, arXiv preprint arXiv:2609.09040 (2026).

\bibitem{LP1998}
F.~Lust-Piquard, \emph{Riesz transforms associated with the number operator on
  the {W}alsh system and the fermions}, Journal of Functional Analysis (1998).

\bibitem{Na2016}
A.~Naor, \emph{Discrete {R}iesz transforms and sharp metric ${X}_p$
  inequalities}, Annals of Mathematics (2016), 991--1016.

\bibitem{N-S2016}
A.~Naor and G.~Schechtman, \emph{Metric {$X_{p}$} inequalities}, Forum Math. Pi
  \textbf{4} (2016), e3, 81. \MR{3456183}

\bibitem{X-Z2026}
Z.~Xu and H.~Zhang, \emph{The endpoint fractional {R}iesz estimate on the
  {H}amming cube}, 2026, Preprint, arXiv:2609.03993.

\end{thebibliography}





\end{document}
